% ========================================================================
% SP1 · material del anexo que pertenece a otros capitulos del TFM
%
% Este fichero NO se compila con sp1.tex. Los anexos son un presupuesto
% compartido por todo el trabajo, de modo que estos bloques se insertan en
% su capitulo de destino al integrar el main. Cada uno conserva el texto
% exacto que tenia.
% ========================================================================

% ------------------------------------------------------------------
% destino: capítulo de metodología

% ------------------------------------------------------------------

\anxsection{Protocolo Monte Carlo pareado}{anx:protocolo}
% --------------------------------------------------------------------------
\subsection{Protocolo Monte Carlo y trazabilidad estadística}

\begin{figure}[H]\centering
\tikzpanel{.98\linewidth}{16.0cm}{figures/protocol_pipeline.tex}
{(b) Diseño $\rightarrow$ registro $\rightarrow$ pareamiento $\rightarrow$ inferencia}
\caption{Protocolo Monte Carlo. Significancia estadística y relevancia práctica se
evalúan por separado cuando ambas condiciones fueron preespecificadas.}
\viuownsource
\end{figure}

% --------------------------------------------------------------------------
% B. Pseudocódigo de los métodos de N3 (procedía de SP1, p. 11)

% ------------------------------------------------------------------
% destino: material de apoyo del repositorio
% ------------------------------------------------------------------


% ------------------------------------------------------------------

\anxsection{Familias decisionales de N4}{anx:familias}
% --------------------------------------------------------------------------
\begin{figure}[H]\centering
\tikzpanelfit{.78\linewidth}{11.2cm}{figures/n4/n4_taxonomy.tex}
{Familias F-I--F-IV}
\caption{Cuatro familias decisionales de N4, la información que comparten y la
salida entera que cada una entrega a SP2.}
\viuownsource
\end{figure}

{\scriptsize
\textbf{Familias de N4: aportación y límite.}\par
\renewcommand{\arraystretch}{1.0}
\begin{tabularx}{\linewidth}{@{}l>{\raggedright\arraybackslash}X>{\raggedright\arraybackslash}Xl@{}}
\toprule
&\textbf{Ventaja}&\textbf{Límite}&\textbf{Ev.}\\
\midrule
F-I&Decide la variable entera que recibe SP2; potencial exacto en desviaciones
unilaterales y terminación demostrada.&Mínimo $h$-local, no global; la etiqueta exige
búsqueda no truncada y aquí se truncó en $h=3$.&N4.E4, N4.E9\\
F-II&Iteración de coste bajo sobre un campo común; compara cuatro dinámicas sin
cambiar el objetivo.&El estado continuo no es una coalición y exige $\mathcal R$;
tras el cierre el orden de mérito cambia en el 74,8\,\%.&N4.E7\\
F-III&Cuotas como restricción compartida, con criterio de aceptación numérico
explícito.&Las KKT motivan el objetivo vGNE sin demostrar convergencia de la
discretización; depende de $\mathcal R$.&N4.E7\\
F-IV&Añade llegadas y fallos en horizonte, más próximo a la operación real.&Banco
reducido y oráculo con información futura: frontera declarada, no contraste
válido.&piloto\\
\bottomrule
\end{tabularx}%
}%

% --------------------------------------------------------------------------
% D. Glosario de la familia F-I (procedía de SP1, p. 21)
% --------------------------------------------------------------------------
{\scriptsize
\textbf{Variantes de F-I.} BR: $h=1$; 2BR: $h=2$; C3: $h=3$; Geo-ASR y Geo-LLL
cambian la regla de revisión, no el orden; CF conserva el orden global de
confirmación; DMIS+TX usa conflictos y versiones. La etiqueta $h$-local solo aplica
si la búsqueda no se trunca.
}

% --------------------------------------------------------------------------
% F. Congelación de configuraciones (procedía de SP1, p. 26)
% --------------------------------------------------------------------------
{\footnotesize Las configuraciones y semillas de SP1 quedaron congeladas antes de la
revisión editorial; guiones y configuraciones se recogen en este anexo.}

% ------------------------------------------------------------------
% G. Diseno de los contrastes de N1 (procedia de SP1, p. 5)

% ------------------------------------------------------------------
% destino: capítulo de metodología
% ------------------------------------------------------------------


% ------------------------------------------------------------------

\anxsection{Diseño de los contrastes de N1}{anx:disenon1}
% ------------------------------------------------------------------
{\scriptsize
\textbf{Tabla 3. Diseño de los contrastes de N1.}\par
\renewcommand{\arraystretch}{1.02}
\begin{tabularx}{\linewidth}{@{}p{1.1cm}p{2.7cm}p{3.15cm}X@{}}
\toprule
&\textbf{Diseño}&\textbf{Estimando}&\textbf{Decisión}\\
\midrule
N1.E1&$5\times3\times3\times100$&ahorro frente a heurística voraz&LCB$_{95}>5\%$ + signo/Holm\\
N1.E4&$2\times5\times30$, 5 perfiles&falso positivo / déficit&McNemar + GEE\\
\bottomrule
\end{tabularx}%
}

Los cuatro experimentos de N1 generaron \NOneRows{} registros.

% ------------------------------------------------------------------
% H. Traza representativa de una ejecucion F-I (procedia de SP1, p. 20)
% ------------------------------------------------------------------
\begin{figure}[H]\centering
\pdfpanel{.90\linewidth}{5.2cm}{assets/figures/n4-v4/n4_fi_representative_trace.pdf}
{Traza F-I}
\caption{$D_4$ y $J_4$ tras cada \emph{commit} aceptado en una ejecución
representativa. Durante la fase de cobertura $J_4$ puede aumentar mientras $D_4$
disminuye.}
\viuownsource
\end{figure}

% ------------------------------------------------------------------
% I. Atlas de reglas y de orden de revision (procedia de SP1, N4)

% ------------------------------------------------------------------
% destino: material de apoyo del repositorio
% ------------------------------------------------------------------


% ------------------------------------------------------------------

\anxsection{Esquema de los comparadores continuos}{anx:continuas}
% ------------------------------------------------------------------
\begin{figure}[H]\centering
\tikzpanel{.82\linewidth}{3.8cm}{figures/n4/n4_population.tex}
{(a) Estado continuo $\rightarrow\mathcal R(x^\star)$}
\par\vspace{.4mm}
\tikzpanel{.82\linewidth}{3.8cm}{figures/n4/n4_primal_dual.tex}
{(b) Primal--dual orientado a vGNE}
\caption{(a) Perfil continuo del símplex y su cierre entero $\mathcal R$.
(b) Intercambio primal-dual con consenso sobre multiplicadores.}
\viuderivedsource{\citet{sandholm2010population}, \citet{quijano2017population},
\citet{yi2019operatorGNE} y \citet{franci2022stochasticGNE}}
\end{figure}


% ------------------------------------------------------------------
% K. Orden conectado frente a mejora (procedia de SP1, N4)

% ------------------------------------------------------------------
% destino: material de apoyo del repositorio
% ------------------------------------------------------------------


% ------------------------------------------------------------------

\anxsection{Qué entrega SP1 a SP2}{anx:interfaz}

SP2 recibe únicamente coaliciones enteras cuya cobertura escalar satisface
el modelo aditivo de SP1, y esa verificación no es una comprobación física. Las
ejecuciones sin coalición válida conservan su estado operacional ---infactible,
divergente, censurada por tiempo o indeterminada--- sin reinterpretarse como
infactibilidad física.

% ------------------------------------------------------------------
% P. Diagnostico post hoc de conectividad (procedia del cuerpo, N3)
% ------------------------------------------------------------------

% ------------------------------------------------------------------
% destino: material de apoyo del repositorio
% ------------------------------------------------------------------


% ------------------------------------------------------------------

\anxsection{Qué responde cada nivel y con qué límite}{anx:preguntas}

{\footnotesize Este resumen se traslada al capítulo de conclusiones, donde se cruzará con las hipótesis globales una vez fijada la arquitectura SP1--SP4 completa.}

{\scriptsize
La evidencia de cada fila son los contrastes N1.E1 y N1.E4, N2.E1 y N2.E2,
N3.E2 y N3.E3, y N4.E4 y N4.E9.\par\smallskip
\renewcommand{\arraystretch}{1.0}
\setlength{\tabcolsep}{4pt}
\begin{tabularx}{\linewidth}{@{}l>{\raggedright\arraybackslash}X
>{\raggedright\arraybackslash}X>{\raggedright\arraybackslash}X@{}}
\toprule
&\textbf{Pregunta}&\textbf{Resultado}&\textbf{Límite}\\
\midrule
N1&¿cuántos AMR hacen falta?&el LP por puestos es integral y la cardinalidad es
exacta bajo homogeneidad&deja de representar la capacidad al crecer el CV\\
N2&¿qué AMR deben formar la coalición?&la composición entera es necesaria y el
óptimo continuo no es ejecutable&certificación sujeta a límite temporal\\
N3&¿qué información necesita una decisión distribuida?&las variantes GRAPE
reducen la brecha a costa de más bytes por AMR&canal fiable; la conectividad
condiciona el éxito\\
N4&¿qué orden de revisión necesita el refinamiento?&$h=2$ recoge la mayor parte
de la mejora por una fracción de los bytes de $h=3$&mínimo $h$-local, búsqueda
truncada en $h=3$\\
\bottomrule
\end{tabularx}
\par\smallskip
El resultado transversal es que, cuando la instancia es factible y el oráculo la
certifica, el método propuesto cubre la demanda sin exceso persistente en los
cuatro niveles del banco descrito; en N1 esa cobertura solo se sostiene bajo
capacidad homogénea. El coste comunicativo se mide siempre frente al oráculo,
pero su escalado con $N$ queda fuera de estas campañas. La correspondencia con
las hipótesis globales del trabajo se establece en el capítulo de conclusiones,
una vez fijada la arquitectura SP1--SP4 completa.
}

% ------------------------------------------------------------------
% R. Resultados de los comparadores continuos (procedia del cuerpo)
% ------------------------------------------------------------------

% ------------------------------------------------------------------
% destino: material de apoyo del repositorio
% ------------------------------------------------------------------


% ------------------------------------------------------------------

\anxsection{Resultados de los comparadores continuos}{anx:continuasres}

La Figura~\ref{fig:continuas-resultados} recoge el detalle que sostiene el
veredicto del cuerpo sobre la vía poblacional.

\begin{figure}[H]\centering
\pdfpanel{.94\linewidth}{6.8cm}{assets/figures/n4-v3/n4_fii_continuous_atomic.pdf}
{(a) N4.E7 · F-II antes y después de $\mathcal R$}
\par\vspace{.8mm}
\pdfpanel{.94\linewidth}{6.8cm}{assets/figures/n4-v3/n4_fiii_residuals.pdf}
{(b) N4.E7 · F-III: residual continuo y coalición cerrada}
\caption{(a) Déficit continuo de las cuatro dinámicas y brecha tras el cierre
$\mathcal R$. (b) Residual del esquema primal-dual y brecha de la coalición cerrada
(N4.E7). Los bytes cubren la fase decisional y excluyen $\mathcal R$.}
\label{fig:continuas-resultados}
\viuownsource
\end{figure}

% ------------------------------------------------------------------
% Mapa de niveles N1--N4 (procedia del cuerpo)
% ------------------------------------------------------------------

% ------------------------------------------------------------------
% destino: material de apoyo del repositorio
% ------------------------------------------------------------------


% ------------------------------------------------------------------

\anxsection{Mapa de niveles N1--N4}{anx:niveles}

Detalle de la Figura~\ref{fig:niveles}.

\begin{figure}[H]\centering
\tikzpanel{.96\linewidth}{16.5cm}{figures/compact/fig02_level_guide.tex}
{Niveles N1--N4}
\caption{Estructura N1--N4 y supuesto que retira cada nivel.}
\label{fig:niveles}
\viuownsource
\end{figure}

% ------------------------------------------------------------------
% Informacion disponible durante la ejecucion (procedia del cuerpo)
% ------------------------------------------------------------------

% ------------------------------------------------------------------
% destino: material de apoyo del repositorio
% ------------------------------------------------------------------


% ------------------------------------------------------------------

\anxsection{Informacion disponible durante la ejecucion}{anx:informacion}

Detalle de la Figura~\ref{fig:n3-informacion}.

\begin{figure}[H]\centering
\tikzpanelfit{.99\linewidth}{6.5cm}{figures/compact/fig14_n3_contract.tex}
{Información disponible durante la ejecución}
\caption{Datos públicos, estado privado de cada AMR y papel del observador durante
una ronda.}
\label{fig:n3-informacion}
\viuownsource
\end{figure}

% ------------------------------------------------------------------
% Barrera, conflictos y confirmacion concurrente (procedia del cuerpo)
% ------------------------------------------------------------------

% ------------------------------------------------------------------
% destino: material de apoyo del repositorio
% ------------------------------------------------------------------


% ------------------------------------------------------------------

\anxsection{Barrera, conflictos y confirmacion concurrente}{anx:concurrencia}

Detalle de la Figura~\ref{fig:concurrencia}.

\begin{figure}[H]\centering
\tikzpanel{.90\linewidth}{5.5cm}{figures/n4/n4_barrier_refinement.tex}
{(a) Barrera y refinamiento}
\par\vspace{.4mm}
\tikzpanel{.90\linewidth}{5.5cm}{figures/n4/n4_conflict_graph.tex}
{(b) Grafo de conflictos}
\par\vspace{.4mm}
\tikzpanel{.90\linewidth}{5.5cm}{figures/n4/n4_tx_flow.tex}
{(c) READ-DMIS-COMMIT}
\caption{(a) Barrera bilateral de \eqref{eq:n4-bilateral-barrier-compact} y
refinamiento progresivo en $h$. (b) Grafo de conflictos entre propuestas con AMR o
carga compartidos. (c) Ciclo READ--DMIS--COMMIT con versiones.}
\label{fig:concurrencia}
\viuderivedsource{la selección distribuida de \citet{luby1986MIS} y el
\emph{commit} versionado de \citet{grayLamport2006Commit}}
\end{figure}

% ------------------------------------------------------------------
% Esquema de CBBA-RB y de las variantes GRAPE (procedia del cuerpo)
% ------------------------------------------------------------------

% ------------------------------------------------------------------
% destino: material de apoyo del repositorio
% ------------------------------------------------------------------


% ------------------------------------------------------------------

\anxsection{Esquema de CBBA-RB y de las variantes GRAPE}{anx:metodos}

Detalle de la Figura~\ref{fig:n3-metodos}.

\begin{figure}[H]\centering
\tikzpanel{.96\linewidth}{4.75cm}{figures/compact/fig15_n3_cbba.tex}
{(a) Capacity-CBBA-RB}
\par\vspace{.5mm}
\tikzpanel{.96\linewidth}{4.75cm}{figures/compact/fig16_n3_grape.tex}
{(b) Weighted-GRAPE / Pair-GRAPE}
\caption{(a) Capacity-CBBA-RB: puja, consenso sobre ganadores y barrera de retorno.
(b) Weighted-GRAPE y Pair-GRAPE: propuesta, inundación del máximo y aplicación
simultánea del perfil común.}
\label{fig:n3-metodos}
\viuderivedsource{\citet{choiBrunetHow2009CBBA}, \citet{johnson2011acbba},
\citet{jangShinTsourdos2018GRAPE} y \citet{dutta2021hedonic}}
\end{figure}

% ------------------------------------------------------------------
% Topologia y robustez del grafo (procedia del cuerpo)
% ------------------------------------------------------------------

% ------------------------------------------------------------------
% destino: material de apoyo del repositorio
% ------------------------------------------------------------------


% ------------------------------------------------------------------

\anxsection{Topologia y robustez del grafo}{anx:topologia}

Detalle de la Figura~\ref{fig:n3-topologia}.

\begin{figure}[H]\centering
\pdfpanel{.98\linewidth}{11.5cm}{assets/figures/sp1-final/f6_n3c_topology_robustness.pdf}
{Topología y robustez}
\caption{(A) Factibilidad condicionada al oráculo por régimen de conectividad. La
partición permanente es control negativo y su barra mide cobertura verificable
\emph{post hoc}, no un perfil global acordado. (B) Bytes/AMR medianos por régimen.
(C) Conectividad algebraica $\lambda_2$ frente a $\Gamma_R$, un punto por mundo y
solo para grafos conexos: con el grafo partido $\lambda_2=0$ y no cabe en el eje
logarítmico (N3.E3).}
\label{fig:n3-topologia}
\viuownsource
\end{figure}

% ------------------------------------------------------------------
% Agregados que intervienen en una desviacion (procedia del cuerpo)

% ------------------------------------------------------------------
% Calidad frente a coste de comunicacion (procedia del cuerpo)
% ------------------------------------------------------------------

% ------------------------------------------------------------------
% destino: material de apoyo del repositorio
% ------------------------------------------------------------------


% ------------------------------------------------------------------

\anxsection{Calidad frente a coste de comunicacion}{anx:pareto}

Detalle de la Figura~\ref{fig:pareto}.

\begin{figure}[H]\centering
\pdfpanel{.94\linewidth}{6.6cm}{assets/figures/n4-v3/n4_cross_family_pareto.pdf}
{Calidad, comunicación y factibilidad}
\caption{Brecha mediana frente a bytes/AMR de los métodos de N3 y N4. El tamaño
del marcador indica la tasa de factibilidad; la forma distingue las salidas
enteras directas de las que requieren el cierre $\mathcal R$, cuyos bytes no se
incluyen. Los puntos proceden de campañas distintas y se presentan como síntesis
descriptiva, no como contraste pareado.}
\label{fig:pareto}
\viuownsource
\end{figure}
